\section{Satuan dan Nilai (px, em, rem, \%, warna)}

Satuan dan nilai menentukan ukuran, jarak, dan tampilan elemen CSS. Pemilihan satuan yang tepat sangat mempengaruhi responsivitas, aksesibilitas, dan kemudahan pemeliharaan stylesheet. Satuan CSS dibagi menjadi dua kategori utama: \textbf{absolut} (ukuran tetap terlepas dari konteks) dan \textbf{relatif} (ukuran yang dihitung berdasarkan elemen lain atau viewport). Pemahaman mendalam tentang setiap satuan membantu pengembang membuat desain yang skalabel dan konsisten di berbagai perangkat \cite{mdnCss}.

\begin{table}[htbp]
\centering
\begin{tabular}{|P{1.5cm}|P{1.5cm}|P{3.5cm}|P{5cm}|}
\hline
\textbf{Satuan} & \textbf{Jenis} & \textbf{Relatif Terhadap} & \textbf{Kapan Digunakan} \\
\hline
\code{px} & Absolut & --- (titik tetap) & Border, shadow; hindari untuk font-size \\
\hline
\code{em} & Relatif & font-size elemen saat ini & Padding/margin relatif ke teks elemen \\
\hline
\code{rem} & Relatif & font-size elemen \code{html} (root) & Font-size, spacing; paling direkomendasikan \\
\hline
\code{\%} & Relatif & Nilai properti parent & Width/height responsif \\
\hline
\code{vw} & Relatif & 1\% lebar viewport & Full-width hero; font-size fluid \\
\hline
\code{vh} & Relatif & 1\% tinggi viewport & Full-height section \\
\hline
\code{ch} & Relatif & Lebar karakter ``0'' & Max-width paragraf (75ch) \\
\hline
\end{tabular}
\caption{Satuan CSS: jenis, basis perhitungan, dan rekomendasi penggunaan}
\end{table}

\textbf{px} (piksel) adalah satuan absolut yang paling umum. Satu px tidak selalu sama dengan satu piksel fisik layar---pada layar HiDPI (Retina), satu CSS px bisa mencakup beberapa piksel fisik. Piksel cocok untuk elemen yang memerlukan ukuran tetap dan presisi, seperti border (\code{border: 1px solid}) dan box-shadow. Namun, penggunaan px untuk \code{font-size} \textbf{tidak disarankan} karena tidak akan berubah saat pengguna memperbesar teks default di pengaturan browser---mengurangi aksesibilitas bagi pengguna dengan gangguan penglihatan \cite{webdevCss}.

\textbf{em} bersifat relatif terhadap \code{font-size} elemen saat ini (atau parent untuk \code{font-size} itu sendiri). Jika \code{font-size: 16px} dan \code{padding: 1.5em}, maka padding menjadi 24px. Kelemahan em: jika elemen bersarang, efek ``compound'' terjadi---em di dalam em dikalikan berulang, membuat ukuran sulit diprediksi. \textbf{rem} (root em) mengatasi masalah ini karena selalu relatif terhadap \code{font-size} elemen \code{<html>} (biasanya 16px). Ini menjadikan rem lebih mudah diprediksi dan direkomendasikan sebagai satuan utama untuk font-size dan spacing \cite{cssTricks}.

Fungsi \code{calc()} memungkinkan perhitungan matematika dengan mencampur satuan yang berbeda. Misalnya, \code{width: calc(100\% - 2rem);} menghitung lebar penuh dikurangi margin. Fungsi ini sangat berguna untuk layout yang memerlukan kombinasi antara satuan relatif dan absolut. Keyword \code{auto} membiarkan browser menghitung nilai terbaik---sering dipakai pada margin untuk centering atau pada width/height untuk ukuran otomatis. Keyword \code{inherit} memaksa properti mengikuti nilai parent, dan \code{initial} me-reset ke nilai bawaan properti \cite{mdnCss}.

\textbf{Contoh:} \code{font-size: 1rem}, \code{width: 25\%}, \code{calc(100\% - 300px)}, \code{max-width: 70ch}.

\begin{konsep}
\textbf{Strategi Satuan yang Direkomendasikan:} Gunakan \code{rem} untuk font-size, \code{rem} atau \code{em} untuk padding/margin, \code{\%} untuk width responsif, \code{px} untuk border/shadow, dan \code{em} untuk media queries. Hindari \code{px} untuk \code{font-size} agar aksesibel \cite{webdevCss}.
\end{konsep}

